% Part: first-order-logic
% Chapter: tableaux
% Section: derivations

\documentclass[../../../include/open-logic-section]{subfiles}

\begin{document}

\iftag{FOL}
      {\olfileid{fol}{tab}{der}}
      {\olfileid{pl}{tab}{der}}

\olsection{\usetoken{P}{tableau}}

\begin{explain}
गृहीतक म्हणजे काय हे आपण सांगितले आहे आणि निगमन नियम दिले आहेत.
!!^{tableau}s यांपासून विगमनाधारित रीतीने निर्माण होतात: प्रत्येक
!!{tableau} ही एकतर एक किंवा अधिक गृहीतकांनी बनलेली एकच शाखा असते, किंवा
!!a{tableau} च्या एखाद्या शाखेवर निगमन नियमांपैकी एक नियम लावल्याने ती मिळते.
\end{explain}

\begin{defn}[!!^{tableau}]
$\sFmla{S_1}{!A_1}$, \dots, $\sFmla{S_n}{!A_n}$ या गृहीतकांसाठीचा
!!^a{tableau} (येथे प्रत्येक $S_i$ हे $\True$ किंवा~$\False$ यांपैकी एक
आहे) म्हणजे पुढील अटी पूर्ण करणारा !!{signed formula}s चा सांत वृक्ष होय:
\begin{enumerate}
\item वृक्षातील सर्वांत वरची $n$ !!{signed formula}s म्हणजे एकाखाली एक
  असलेली $\sFmla{S_i}{!A_i}$ ही सूत्रे होत.
\item वृक्षातील गृहीतक नसलेली प्रत्येक !!{signed formula} तिच्या वरच्या
  शाखेतील !!a{signed formula} ला निगमन नियम योग्य रीतीने लावल्याने मिळते.
\end{enumerate}
!!a{tableau} ची एखादी शाखा $\sFmla{\True}{!A}$ आणि
~$\sFmla{\False}{!A}$ ही दोन्ही सूत्रे धारण करत असेल तर आणि तरच ती
\emph{बंद} असते; अन्यथा ती \emph{खुली} असते. ज्याची प्रत्येक शाखा बंद
असते तो !!^a{tableau} (त्याच्या गृहीतकांच्या संचासाठीचा)
\emph{बंद !!{tableau}} असतो. !!a{tableau} बंद नसेल, म्हणजे त्यात किमान
एक खुली शाखा असेल, तर तो \emph{खुला} असतो.
\end{defn}

\begin{ex}
  गृहीतकांचा प्रत्येक संच स्वतःच !!a{tableau} असतो, परंतु साधारणतः तो बंद
  नसतो. (स्पष्टपणे, तो बंद असेल तर गृहीतकांत $\sFmla{\True}{!A}$ आणि
  ~$\sFmla{\False}{!A}$ ही !!{signed formula}s ची जोडी आधीच असली पाहिजे.)

  एखाद्या !!a{tableau} मधील !!a{signed formula}~$!A$ ला निगमन नियमांपैकी
  एक नियम लावून त्या टॅब्लोपासून (तो खुला असो वा बंद) नवा आणि अधिक मोठा
  टॅब्लो मिळवता येतो. नियम लागू केलेली~$!A$ ची आस्थिती ज्या ज्या शाखेत
  आहे, त्या प्रत्येक शाखेच्या शेवटी तो नियम एक किंवा अधिक
  !!{signed formula}s जोडतो.

  उदाहरणार्थ, $\sFmla{\True}{!A \land \lnot !A}$ हे गृहीतक विचारात घ्या.
  केवळ त्या गृहीतकाने बनलेला (खुला) !!{tableau} पुढीलप्रमाणे आहे:
  \begin{center}
    \begin{tableau}{}
      [\sFmla{\True}{\formula{A} \land \lnot \formula{A}}, just=\TAss]
    \end{tableau}{}
  \end{center}
  त्या गृहीतकाला $\TRule{\True}{\land}$ नियम लावून त्यापासून नवा
  !!{tableau} मिळतो. हा नियम !!{tableau} मध्ये $\sFmla{\True}{!A}$ आणि
  $\sFmla{\True}{\lnot !A}$ या दोन नव्या ओळी जोडू देतो:
  \begin{center}
    \begin{tableau}{}
      [\sFmla{\True}{\formula{A} \land \lnot \formula{A}}, just=\TAss
        [\sFmla{\True}{\formula{A}}, just={\TRule{\True}{\land}[1]},
          [\sFmla{\True}{\lnot \formula{A}}, just={\TRule{\True}{\land}[1]}
          ]
        ]
      ]
    \end{tableau}{}
  \end{center}
  !!{tableau}s लिहिताना आपण लावलेले नियम उजवीकडे नोंदवतो (उदा.,
  $\TRule{\True}{\land} 1$ याचा अर्थ त्या ओळीवरील !!{signed formula} ही
  ओळ~$1$ वरील !!{signed formula} ला $\TRule{\True}{\land}$ नियम लावल्याने
  मिळाली आहे). या नव्या !!{tableau} मध्ये आता अधिक !!{signed formula}s
  आहेत, परंतु त्यांपैकी फक्त एकीला ($\sFmla{\True}{\lnot !A}$) नियम लावता
  येतो (या स्थितीत $\TRule{\True}{\lnot}$ नियम). त्यामुळे पुढील बंद
  !!{tableau} मिळतो:
  \begin{center}
    \begin{tableau}{}
      [\sFmla{\True}{\formula{A} \land \lnot \formula{A}}, just=\TAss
        [\sFmla{\True}{\formula{A}}, just={\TRule{\True}{\land}[1]}
          [\sFmla{\True}{\lnot \formula{A}}, just={\TRule{\True}{\land}[1]}
            [\sFmla{\False}{\formula{A}}, just={\TRule{\True}{\lnot}[3]}, close]
          ]
        ]
      ]
    \end{tableau}{}
  \end{center}
\end{ex}

\end{document}
